%%%%%%%%%%%%%%%%%%%%%%%%
% $Autor: MansourAhmadyPhoulady, RanjeetsingTawar, HemanthPrabhakaran$
% $Pfad: Car_Licence_Plate_Identification_with_a_JetsonNano$
% $Version: 4.0.0 $
%
% !TeX encoding = utf8
%
%%%%%%%%%%%%%%%%%%%%%%%%

\chapter{Development Folder Structures}

\section{Folder Structure}

The project directory is meticulously organized to ensure efficient navigation, development, and debugging. Below is the detailed structure of the folder hierarchy and its contents:

\begin{itemize}
	\item \texttt{.ipynb\_checkpoints/}: Contains auto-saved checkpoints of Jupyter notebooks, providing a backup mechanism for ongoing development and version control.
	\item \texttt{anprsys/}: This directory is the core of the project, housing the main application scripts, pre-trained TensorFlow models, configuration files, and essential environment setup scripts.
	\item \texttt{Tensorflow/}: Dedicated to TensorFlow-specific resources, including model configurations, training scripts, and auxiliary files required for the machine learning tasks.
	\item \texttt{.gitignore}: A configuration file specifying which files and directories Git should ignore, such as temporary files, logs, and sensitive data.
	\item \texttt{README.md}: The primary documentation file providing a comprehensive overview of the project, including its purpose, setup instructions, and usage guidelines.
	\item \texttt{report.txt}: Stores generated logs and reports from the system's operation, capturing key performance metrics and results for analysis.
	\item \texttt{Error Guide.md}: A markdown file detailing common issues and their troubleshooting steps to assist developers in debugging efficiently.
	\item \textbf{Jupyter Notebooks:}
	\begin{itemize}
		\item \texttt{Automatic Number Plate Detection.ipynb}: A fully interactive notebook for executing the number plate detection pipeline, including data preprocessing, training, and testing.
	\end{itemize}
	\item \texttt{Detection\_Images/}: A repository of output images where detected license plates are highlighted, serving as visual proof of successful detections.
	\item \texttt{detection\_results.csv}: A structured CSV file storing validated license plate detections along with their corresponding UUIDs for reference and traceability.
	\item \texttt{realtimeresults.csv}: A real-time log of license plates detected using a live feed from the connected Logitech webcam, stored alongside UUIDs for accuracy tracking.
\end{itemize}

\section{Overview}

For the project to function seamlessly, the following prerequisites are essential:

\begin{itemize}
	\item A Python programming environment configured for machine learning and image processing tasks.
	\item TensorFlow framework for handling deep learning processes.
	\item Supporting libraries such as OpenCV for image processing and easyOCR for optical character recognition.
	\item A \texttt{requirements.txt} file detailing the exact versions of the libraries and packages required for consistent performance across development environments.
\end{itemize}

\section{Installation}

To set up the environment, follow these steps:

\begin{enumerate}
	\item \textbf{Create a Virtual Environment:} 
	\begin{verbatim}
		python -m venv venv
		source venv/bin/activate  # For Linux/MacOS
		venv\Scripts\activate     # For Windows
	\end{verbatim}
	\item \textbf{Install Dependencies:} Use the provided \texttt{requirements.txt} file to install all necessary libraries and dependencies:
	\begin{verbatim}
		pip install -r requirements.txt
	\end{verbatim}
	\item \textbf{Configure TensorFlow Models:} Obtain pre-trained models from the TensorFlow Model Zoo or official TensorFlow repositories and configure them for object detection tasks.
\end{enumerate}

\section{First Steps}

After completing the installation, begin the development workflow as follows:

\begin{enumerate}
	\item \textbf{Activate the Virtual Environment:} Ensure all dependencies are active by starting the virtual environment.
	\item \textbf{Utilize Jupyter Notebooks:} Open the interactive notebooks provided:
	\begin{itemize}
		\item Start with \texttt{1. Image Collection.ipynb} to gather and preprocess training data.
		\item Proceed to \texttt{2. Training and Detection.ipynb} for model training and initial testing.
	\end{itemize}
\end{enumerate}

\section{Example}

The following sequence demonstrates a typical use-case workflow of the project:

\begin{enumerate}
	\item \textbf{Step 1 - Data Collection:} Execute the notebook \texttt{1. Image Collection.ipynb} to collect images for training, capturing diverse examples of license plates under varying conditions.
	\item \textbf{Step 2 - Model Training:} Run \texttt{2. Training and Detection.ipynb} to train the detection model, leveraging TensorFlow’s pre-trained capabilities and fine-tuning it with the collected data.
	\item \textbf{Step 3 - Real-Time Detection:} Deploy the trained model on live data to detect and recognize license plates, storing the results in the \texttt{detection\_results.csv} and \texttt{realtimeresults.csv} files for review.
	\item \textbf{Step 4 - Output Analysis:} Analyze the stored results and detection images in the \texttt{Detection\_Images/} directory to evaluate system performance and accuracy.
\end{enumerate}

\section{Conclusion}

By structuring the development folder comprehensively and following the outlined steps, the project ensures an efficient workflow. The organized structure, combined with clear installation and operational guidelines, supports seamless development, testing, and deployment, making it easier to maintain and scale the system over time.
